\documentclass[11pt,a4paper]{article}
\usepackage[T1]{fontenc}
\usepackage[utf8]{inputenc}
\usepackage{lmodern}
\usepackage[a4paper,margin=1in]{geometry}
\usepackage{microtype}
\usepackage{enumitem}
\usepackage{amsmath,amssymb}
\usepackage{mathtools}
\usepackage{hyperref}
\hypersetup{
	colorlinks=true,
	linkcolor=black,
	urlcolor=blue,
	citecolor=black
}

\setlist[itemize]{leftmargin=1.2em}
\setlist[enumerate]{leftmargin=1.2em}

\newcommand{\Title}[1]{\section*{#1}\vspace{-0.25em}\hrule\vspace{0.8em}}
\newcommand{\Field}[2]{\textbf{#1:}~#2\par}
\newcommand{\Affil}[1]{\textit{#1}\par}
\newcommand{\Abstract}{\vspace{0.6em}\textbf{Abstract.}~}

\begin{document}
	
	\begin{center}
		{\LARGE Online Algebra, Logic, and Topology Seminar}\\[0.25em]
		{\large Category Theory Seminar — Abstract}\\[0.5em]
		October 07, 2025
	\end{center}
	
	\vspace{1.25em}\hrule\vspace{1.25em}
	
	\Title{When are exact categories co-exact?}
	
	\Field{Speaker}{James Gray}
	\Affil{Stellenbosch University, South Africa}
	
	\Field{Date \& Time}{October 07, 2025 (Tuesday) — 3{:}00\,PM (Coimbra time)}
	\Field{Place}{Google Meet}
	
	\Abstract
	The aim of the talk is two-fold: (i) to explain that finitely co-complete
	pretoposes are co-ideally-exact~\cite{J} and co-arithmetical~\cite{P}, and (ii)
	to explain that there are conditions in common between additive and
	lextensive categories, which together with exactness (and finite
	cocompleteness) imply co-exactness and co-protomodularity. As
	a consequence of (i) we recover Theorem~1.3 of~\cite{BC} which states that
	the dual of the category of pointed compact Hausdorff spaces is
	semi-abelian, which in fact was the motivation to begin this work.
	
	\begin{thebibliography}{9}
		\bibitem[BC]{BC} F.~Borceux, M.M.~Clementino, \emph{On toposes, algebraic theories,
			semi-abelian categories and compact Hausdorff spaces}, Theory and
		Applications of Categories 43(11), 363--381, 2025.
		
		\bibitem[J]{J} G.~Janelidze, \emph{Ideally exact categories}, Theory and Applications of
		Categories 41(11), 414--425, 2024.
		
		\bibitem[P]{P} M.C.~Pedicchio, \emph{A categorical approach to commutator theory},
		Journal of Algebra 177, 647--657, 1995.
	\end{thebibliography}
	
\end{document}
